\section{Learning the cost of a complete reduction}
\label{app:workflow}
This appendix defines the complete-time prediction problem, derives the fitted models, and specifies the physical inputs used in Section~\ref{sec:workflow}. The cost model receives a feature vector $\boldsymbol f_c$ for target batch $c$.
Its entries are the number of loops, external momenta, inverse propagators,
physical propagators and symbolic invariants; $\log(1+n_c)$, where $n_c$ is
the number of targets; the maximum dot count, maximum numerator rank, mean dot count and mean numerator rank;
the number of target sectors; and seven indicators for the graph structures.
An indicator is one for its graph and zero for the other graphs.
Together these quantities give eighteen inputs, all available from the requested targets and family before reduction.

For an index vector $\boldsymbol a$, define
\begin{equation}
 r(\boldsymbol a)=\sum_i\max(a_i,0),\qquad
 s(\boldsymbol a)=\sum_i\max(-a_i,0),\qquad
 q(\boldsymbol a)=\sum_i\max(a_i-1,0).
\end{equation}
Here $r$ is the positive index sum, $s$ is the numerator rank, and $q$ is
the number of extra denominator powers, called dots. These definitions
follow directly by counting the denominator and numerator factors.
The initial bounds $B_c=(r_c,s_c,q_c)$ are the componentwise maxima over
the targets. The trial seed regions use $B_c$, $B_c-(1,1,1)$,
$B_c-(1,0,1)$, or $B_c+(1,1,1)$. Negative bounds are set to zero.
The choices change the region in which equations are generated. Since a smaller region can omit required relations, every attempt must still pass the prescribed-basis check.

Let $t_{cj}$ be the median complete time of strategy $j$ on training batch
$c$, and let $t_{c0}$ be the time of the integrated native strategy.
The training response is
\begin{equation}
 y_{cj}=\log\frac{t_{cj}}{t_{c0}}.
 \label{eq:workflow-label}
\end{equation}
Each time includes unsuccessful attempts and the eventual successful fallback, so a strategy that omits necessary equations is penalized by the extra calculation. The logarithm preserves the ordering of positive time ratios; minimizing the response therefore also minimizes the complete time relative to the common baseline.

We fit constant, tree, linear and nearest-neighbour predictors of this
response. The constant predictor averages each strategy's training
responses. A tree splits the feature space and uses a constant vector
within each leaf, following regression-tree construction~\cite{CART1984}.
For a leaf containing training set $A$, minimizing
$\sum_{c\in A}(y_{cj}-v_j)^2$ gives
\begin{equation}
 \frac{\partial}{\partial v_j}
 \sum_{c\in A}(y_{cj}-v_j)^2
 =2|A|v_j-2\sum_{c\in A}y_{cj}=0,
 \qquad v_j=\frac{1}{|A|}\sum_{c\in A}y_{cj}.
\end{equation}
Each split minimizes the sum of squared errors over the seven strategies.
The tested maximum depths are two, four and six. Each depth is tested
with a minimum of either one or three training batches per child.

For the linear models, we first subtract the training mean of each feature and
divide by its training population standard deviation, computed with
divisor equal to the number of training batches. A constant feature uses
divisor one. Let $Z$ contain these standardized feature rows with a leading intercept column, and let $Y$ contain the responses.
Ridge regression minimizes
\begin{equation}
 J(W)=\|ZW-Y\|_F^2+\lambda\operatorname{tr}(W^TPW),
 \qquad P=\operatorname{diag}(0,1,\ldots,1).
\end{equation}
Here $W$ is the matrix of fitted coefficients, $\|\cdot\|_F^2$ is the sum
of squared matrix entries, and $P$ leaves the intercept unpenalized.
Differentiating each entry of $W$ gives
\begin{equation}
 \nabla_WJ=2Z^T(ZW-Y)+2\lambda PW=0,
 \qquad (Z^TZ+\lambda P)W=Z^TY.
\end{equation}
We solve this system for $\lambda=0.1,1,10$.

The nearest-neighbour models estimate a cost from similar training batches. Eligible neighbours must have the same graph, loop count, external-momentum count, propagator counts, and number of symbolic invariants. They choose up to two or four training batches with
the smallest squared distance in standardized features. The predicted
response is the mean of their response vectors. The leaf-mean derivative
above also derives this local constant estimate. If no eligible neighbour
exists, the predictor returns equal costs and the selector uses native
choices.

Let $j_*$ be the strategy with the smallest predicted logarithmic cost
$\widehat y_j(\boldsymbol f)$. Validation selects a margin
$\gamma\in\{0,0.05,0.15\}$, and the model chooses the strategy as follows:
\begin{equation}
 \widehat j(\boldsymbol f)=
 \begin{cases}
 j_*,&\widehat y_{j_*}(\boldsymbol f)<\widehat y_0(\boldsymbol f)-\gamma,\\
 0,&\text{otherwise}.
 \end{cases}
 \label{eq:workflow-choice}
\end{equation}
Subtracting the two predicted logarithmic costs shows that the condition
requires a predicted relative cost below $e^{-\gamma}$. Validation selects
the model and margin with the smallest geometric mean time ratio.
The selected settings remain fixed during reserved tests.
This finite comparison establishes the best measured candidate within
the study. It does not prove a globally optimal reduction strategy.

\subsection{Parameter points and target lists}
The comparison uses the sunrise, planar box, three-loop vacuum, three-loop ladder, nonplanar box, pentagon, and box with external virtuality, with the same denominator routings and kinematic conventions as the preceding tests. The definitions below specify the parameter values and target lists for each split. Let $h$ identify a rational
parameter profile. Training uses $h=41,42$, validation uses $h=43$, and
testing uses $h=44,45$. In a common squared-momentum unit, the sunrise has
squared masses $(2,7,11)$ and $p^2=-1009-12h$. The planar box has
$(s,t)=(-1009-12h,-1021-14h)$. The three-loop vacuum has squared masses
$(1,1,1,1,1,53+2h)$. The ladder keeps squared masses
$(1,0,0,0,0,0,0,0)$ and has $p^2=-1009-12h$.
The scattering assignments are
\begin{align}
 (s,t)_{\rm NP}&=(-113-18h,-127-10h),\\
 (s_{12},s_{23},s_{34},s_{45},s_{51})_{\rm P}
 &=(-11-20h,-13-20h,-17-20h,-19-20h,-23-20h),\\
 (s,t,v)_{\rm box}&=(-113-18h,-127-10h,-149-12h).
\end{align}
These assignments fix the rational parameter values in each profile while leaving the dimension symbolic. The additional symbolic batches instead retain all independent invariants of their scattering graph as symbols.

For the four earlier multiloop graphs, we retain the ordered target list
that starts with the corner, raises each physical denominator power once,
and then introduces each single auxiliary numerator. We repeat these
steps after removing one, two and three physical denominators, retain at
least two, and remove duplicates. Each rational profile supplies a
four-target batch and a batch of the first 32 targets, or all 21 targets
for the sunrise. The three scattering graphs use the same construction
with at most one removed denominator. Their larger batches have 32, 31
and 21 targets for the nonplanar box, pentagon and virtuality box.

For the symbolic nonplanar box and pentagon, training takes target entries
12--15, validation takes 16--19, and testing takes 20--23, with numbering
starting at zero. These lists do not overlap. For the symbolic virtuality
box, the respective single targets are $(3,1,1,1)$, $(1,3,1,1)$ and
$(1,1,3,1)$. These targets have two extra powers on one physical denominator.
The rational batches reuse target patterns at different parameter points.

The joint-request test uses the larger rational batch at $h=44$ for each
graph and the three symbolic test batches. The symbolic virtuality box
also includes $(1,1,1,3)$, so it can supply two disjoint requests. A group
of $n$ targets is divided after entry $\lfloor n/2\rfloor$, with the first
request containing that many entries. The two requests share the same family definition and are available together, and their union contains each target once. By linearity, selecting the corresponding target rows from the joint reduction returns the coefficients for each original request. This separation leaves every coefficient unchanged and requires no new integral identity. The test reuses reserved parameter points and keeps the learned
model fixed.

\subsection{Combining the measured times}
Let $m_{cj}$ be the median of the five test times for batch $c$ and
method $j$, and let method zero be its stated native control.
For graph $g$ with batch set $C_g$, we define
\begin{equation}
 R_{gj}=\exp\!\left[\frac{1}{|C_g|}
       \sum_{c\in C_g}\log\frac{m_{cj}}{m_{c0}}\right],
 \qquad
 R_j=\exp\!\left[\frac{1}{7}\sum_{g=1}^{7}\log R_{gj}\right].
\end{equation}
The logarithm of a product equals the sum of the logarithms, so each
expression is the corresponding geometric mean of the time ratios.
The second expression assigns equal weight to each graph even when the
numbers of batches differ. A ratio $R_j<1$ means less time. The percentage
saving is $100(1-R_j)$, and the reciprocal $1/R_j$ is the speedup.

For the descriptive intervals, we sample seven graph labels with
replacement, calculate their mean logarithmic ratio, and exponentiate.
We repeat this 10,000 times and report the 2.5th and 97.5th percentiles.
This procedure describes variation among the seven tested graphs.
It does not provide a population guarantee for other integral families.
The joint-request comparison uses the same calculation, with separate
native requests as its first control. Its parameter points overlap the
main reserved test, so it is an additional comparison on the same test
population.
